\customSection[sec:analysis]{Аналіз предметної області та інженерія вимог вебзастосунку}{0em}

\customSubSection[sec:problem]{Проблематика планування електромережі приміщень і розрахунку матеріалів}{1em}

\customSubSubSection[sec:problem-wiring]{Розміщення електрообладнання та трасування проводки}{0.5em}

Внутрішньобудинкова електромережа приміщення складається з електрощитка, розподільних коробок, групових ліній, розеток, вимикачів і світильників; нормативні вимоги до її проєктування в житлових будівлях встановлює ДБН~В.2.5-23~\citecustom{6}. На етапі планування, що передує складанню проєктної документації, виникає прикладна задача: визначити, на якій стіні й на якій висоті розташувати кожен елемент, як провести між ними проводи та скільки кабелю й оздоблювальних матеріалів для цього потрібно.

Розташування елемента на стіні однозначно задається чотирма параметрами: кімнатою, стіною, відстанню від кута стіни та висотою над підлогою. Так само визначаються двері й арки, світильники та предмети на стінах. Така «стінова» система координат природна для виконавця робіт, оскільки відповідає тому, як розмітку виконують на об'єкті: від кута та від підлоги.

Довжина проводу залежить від маршруту, а не від відстані між елементами за прямою. У прийнятій у застосунку моделі провід прокладають уздовж стін і стелі, а повороти виконують під прямим кутом; у такій моделі довжина дорівнює сумі довжин ортогональних відрізків (манхеттенська відстань), а якщо елементи розташовані на різних стінах або в різних кімнатах, маршрут проходить через стелю. На контрольній сцені, описаній у підрозділі~\ref{sec:verification}, прямолінійна відстань від щитка до трьохгніздої розетки становить приблизно 3,00~м, тоді як ортогональний маршрут через стелю має довжину 7,10~м, тобто більш ніж удвічі більшу. Оцінка довжини «на око» за прямою, таким чином, систематично занижує потребу в кабелі.

\customSubSubSection[sec:problem-materials]{Розрахунок оздоблювальних матеріалів}{0.5em}

Після розміщення елементів і прокладання проводів потрібно підготувати специфікацію матеріалів: листи ГКЛ для обшивки стін і стелі, профілі каркаса, плінтуси та кабель. Базовий підхід до кількості листів зводиться до формули: площа поверхні з урахуванням запасу ділиться на площу листа та округлюється вгору. Онлайн-калькулятори гіпсокартону працюють саме так: за введеною площею, розміром листа й відсотком запасу вони повертають кількість листів, при цьому прорізи або віднімаються від площі~\citecustom{11}, або задаються кількістю дверей і вікон із фіксованою площею~\citecustom{10}. Розкладки листів ці засоби не будують.

Простий розрахунок за площею має три недоліки. По-перше, запас у відсотках не пов'язаний із реальною геометрією: той самий відсоток дає різну кількість зайвих листів для різних кімнат. По-друге, відкоси дверей і арок, які теж обшивають, у простих калькуляторах не враховуються. По-третє, площа не враховує, що лист не можна розрізати довільно: фрагменти від розрізу бувають придатними для повторного використання, а можуть бути нерозривними відходами.

\customSubSubSection[sec:problem-cutting]{Задача розкрою листового матеріалу}{0.5em}

Задача визначення мінімальної кількості листів для покриття набору прямокутних поверхонь належить до класу задач розкрою та пакування. Її класифікація за структурою вхідних даних і цільовою функцією наведена в уніфікованій типології~\citecustom{1}; класична задача одновимірного розкрою (cutting stock problem) розв'язується із застосуванням лінійного програмування~\citecustom{2}, а огляд двовимірних задач пакування й розкрою наведено в~\citecustom{3}. Для двовимірного випадку важливий клас гільйотинних розкроїв, у яких кожен розріз проходить від краю до краю прямокутника; саме так працює різання ГКЛ різаком, і для таких розкроїв розроблено точні й евристичні алгоритми~\citecustom{4}. Двовимірні задачі пакування узагальнюють задачу про упаковку в контейнери, яка є NP-складною, тому на практиці широко використовують евристики, зокрема жадібні схеми з упорядкуванням елементів за спаданням розміру, оцінки якості яких досліджено в~\citecustom{5}.

У межах цієї роботи задачу сформульовано так. Задано множину прямокутних поверхонь $S_k$ розмірами $W_k\times H_k$ (стіни, стеля, відкоси), множину прямокутних прорізів $O_{kj}\subset S_k$, які покриття не потребують, і розмір стандартного листа $s_w\times s_h$. Потрібно покрити чисту площу $\bigcup_k\bigl(S_k\setminus\bigcup_j O_{kj}\bigr)$ шматками, вирізаними гільйотинними розрізами зі стандартних листів і з відрізків, які залишилися, мінімізуючи кількість використаних нових листів $N$. Додаткові вимоги: повторне використання відрізків, зокрема вирізок із місць дверей і арок, а також (як м'яка вимога) шахове зміщення швів. Очевидною нижньою межею є
\begin{equation}\label{eq:lowerbound}
    N \;\ge\; \left\lceil \frac{A_{\mathrm{net}}}{s_w\, s_h} \right\rceil ,
\end{equation}
де $A_{\mathrm{net}}$~-- сумарна чиста площа поверхонь. Цією межею користуються в підрозділі~\ref{sec:verification} для оцінки якості розкладки.

\customSubSubSection[sec:problem-statement]{Постановка проблеми}{0.5em}

Аналіз дозволяє виділити такі протиріччя, які визначають завдання на розроблення:
\begin{itemize}[noitemsep, topsep=0.3em, leftmargin=*]
    \item між просторовою природою задачі (стіни, висоти, маршрути) та плоским характером ручних підрахунків і електронних таблиць;
    \item між потребою знати довжину проводу за реальним маршрутом і відсутністю маршруту як об'єкта в засобах планування інтер'єру;
    \item між простою формулою кількості листів і реальною розкладкою з розрізами, відкосами та відрізками;
    \item між потребою працювати з проєктом без встановлення програм і складністю десктопних засобів.
\end{itemize}

Таким чином, \textbf{проблема} формулюється як розроблення вебзастосунку, який у браузері дозволяє змоделювати приміщення з електрообладнанням і проводкою, автоматично визначає довжини проводів за ортогональними маршрутами та специфікацію оздоблювальних матеріалів і будує розкладку листів ГКЛ із мінімізацією відходів.

\customSubSection[sec:analogs]{Огляд програмних аналогів}{1em}

Для визначення ніші розглянуто три групи засобів. Перша~-- редактори планування приміщень. Sweet Home 3D~-- вільна програма для інтер'єрного дизайну з побудовою стін, вставкою дверей і вікон та синхронним 3D-переглядом; працює у вигляді настільних застосунків і онлайн-версії~\citecustom{7}. Планувальник Planner~5D є браузерним 2D/3D-редактором із каталогом понад 5000 об'єктів~\citecustom{9}. RoomSketcher дозволяє креслити електричні плани, використовуючи бібліотеку понад 250 символів (розетки, вимикачі, світильники тощо), але повний набір електричних символів доступний лише в тарифах Pro і Team, а проводка в ньому є символом, а не маршрутом із довжиною~\citecustom{8}.

Друга група~-- онлайн-калькулятори гіпсокартону. Калькулятор~\citecustom{11} віднімає площу прорізів, додає відсоток запасу й повертає кількість листів; калькулятор~\citecustom{10} приймає площу стін, кількість дверей і вікон із фіксованою площею (2 і 1,5~м$^2$ відповідно), товщину та розмір листа й запас. Жоден із них не будує розкладку листів і не згадує відкоси.

Третя група~-- поєднання цих можливостей в одному застосунку. Зіставлення за ключовими критеріями наведено в \tabref{tab:analogs}. Значення «не заявлено» означає, що відповідну можливість не згадано в описі функцій на сторінці виробника на дату перегляду.

\begin{spacing}{1.15}
\small\noindent
\begin{xltabular}{\textwidth}{| >{\raggedright\arraybackslash}p{2.9cm} | >{\raggedright\arraybackslash}X | >{\raggedright\arraybackslash}X | >{\raggedright\arraybackslash}X | >{\raggedright\arraybackslash}X |}
    \caption{Порівняння засобів планування та розрахунку матеріалів} \label{tab:analogs} \\
    \hline
    \textbf{Засіб} & \textbf{Формат і 3D} & \textbf{Електрообладнання та проводка} & \textbf{Розрахунок матеріалів} & \textbf{Розкладка листів} \\ \hline
    \endfirsthead

    \hline
    \textbf{Засіб} & \textbf{Формат і 3D} & \textbf{Електрообладнання та проводка} & \textbf{Розрахунок матеріалів} & \textbf{Розкладка листів} \\ \hline
    \endhead

    Sweet Home 3D~\citecustom{7} & настільний і онлайн; 3D-перегляд & не заявлено & не заявлено & не заявлено \\ \hline
    Planner 5D~\citecustom{9} & браузер; 2D і 3D & не заявлено & не заявлено & не заявлено \\ \hline
    RoomSketcher~\citecustom{8} & браузер; плани приміщень & символи (у платних тарифах), проводка як символ & не заявлено & не заявлено \\ \hline
    Калькулятори ГКЛ~\citecustom{10}, \citecustom{11} & вебсторінка; без моделі & ні & кількість листів за площею та запасом & ні \\ \hline
    \textbf{Розроблюваний застосунок} & браузер; 3D-модель кількох кімнат & 10 типів елементів, світильники, проводи з ортогональними маршрутами & площі, ГКЛ, профілі, плінтуси, довжина проводів & так: 4 режими, відрізки, нумерація \\ \hline
\end{xltabular}
\end{spacing}

Огляд дозволяє зробити три висновки. По-перше, засоби планування інтер'єру не пов'язують розміщення електрообладнання з довжинами проводів і витратами матеріалів. По-друге, калькулятори матеріалів не мають просторової моделі, тому не можуть врахувати реальну геометрію прорізів і відкосів та побудувати розкладку. По-третє, поєднання тривимірної моделі, трасування проводів і розкрою листів в одному безкоштовному браузерному застосунку не заявлено в жодному з розглянутих рішень. Ці висновки отримано з відкритих описів виробників і мають розглядатися як гіпотеза, що потребує уточнення в разі зміни функціоналу аналогів.

\customSubSection[sec:requirements]{Специфікація вимог до вебзастосунку}{1em}

\customSubSubSection[sec:stakeholders]{Зацікавлені сторони}{0.5em}

Вимоги сформульовано за рекомендаціями ISO/IEC/IEEE 29148:2018~\citecustom{12}, який визначає процеси інженерії вимог до систем і програмних продуктів та формат їх запису; додатково використано підходи до специфікування вимог і прецедентів із~\citecustom{14}, \citecustom{15}. Визначено чотири групи зацікавлених сторін (\tabref{tab:stake}).

\begin{spacing}{1.15}
\small\noindent
\begin{xltabular}{\textwidth}{| >{\raggedright\arraybackslash}p{4.2cm} | >{\raggedright\arraybackslash}X | >{\raggedright\arraybackslash}X |}
    \caption{Зацікавлені сторони та їхні інтереси} \label{tab:stake} \\
    \hline
    \textbf{Сторона} & \textbf{Інтерес} & \textbf{Основні потреби} \\ \hline
    \endfirsthead

    \hline
    \textbf{Сторона} & \textbf{Інтерес} & \textbf{Основні потреби} \\ \hline
    \endhead

    Проєктувальник / виконроб & швидко спланувати розміщення електрообладнання та підготувати специфікацію & точні розміри, наочна 3D-модель, довжини проводів, кількість матеріалів \\ \hline
    Електромонтажник & знати розташування елементів і маршрути & відстані від кутів і висоти від підлоги, маршрути з контрольними точками \\ \hline
    Власник житла & побачити результат до початку робіт і оцінити витрати & проста робота в браузері, наочність, перелік матеріалів \\ \hline
    Розробник / супровід & розвивати застосунок без складної інфраструктури & модульна структура, відсутність збірника та сервера, відтворювані розрахунки \\ \hline
\end{xltabular}
\end{spacing}

\customSubSubSection[sec:fr]{Функціональні вимоги}{0.5em}

Функціональні вимоги подано як перевірні твердження з критерієм приймання (\tabref{tab:fr}). Результати перевірки вимог наведено в підрозділі~\ref{sec:verification}.

\begin{spacing}{1.15}
\small\noindent
\begin{xltabular}{\textwidth}{| c | >{\raggedright\arraybackslash}X | >{\raggedright\arraybackslash}X |}
    \caption{Функціональні вимоги} \label{tab:fr} \\
    \hline
    \textbf{ID} & \textbf{Вимога} & \textbf{Критерій приймання} \\ \hline
    \endfirsthead

    \hline
    \textbf{ID} & \textbf{Вимога} & \textbf{Критерій приймання} \\ \hline
    \endhead

    FR-01 & Система дозволяє створити новий проєкт або імпортувати наявний; до цього інші вкладки недоступні & вкладки «Конструктор», «Проводка», «Розрахунки» заблоковані до створення чи імпорту \\ \hline
    FR-02 & Система дозволяє задати кілька кімнат із розмірами, позицією, рівнем підлоги та товщиною кожної стіни & кімнати відображаються у 3D, спільні стіни можна сумістити \\ \hline
    FR-03 & Система дозволяє додати двері та арки з параметрами (стіна, зсув від кута, розміри, відкоси, поріг, лиштва, заокруглення, відчинення) & об'єкт з'являється на стіні, прорізи враховуються в розрахунках \\ \hline
    FR-04 & Система дозволяє розмістити електрообладнання (10 типів), світильники (стельові й настінні) та предмети на стінах & елемент розташований у заданих відстані від кута та висоті \\ \hline
    FR-05 & Система дозволяє інтерактивно переміщувати об'єкти у 3D, зокрема між стінами та кімнатами & після переміщення зсув і висота оновлюються у формі \\ \hline
    FR-06 & Система дозволяє прокласти провід між двома вузлами з ортогональним маршрутом, контрольними точками, типом кабелю, кількістю жил і діаметром & усі повороти маршруту прямі; точки можна додавати, тягнути й видаляти \\ \hline
    FR-07 & Система розраховує площі стін кожної кімнати з відніманням прорізів & збіг з ручним розрахунком до 0,01~м$^2$ \\ \hline
    FR-08 & Система розраховує кількість листів ГКЛ за площею з урахуванням відкосів, стелі та відсотка відходів & збіг з ручним розрахунком у штуках \\ \hline
    FR-09 & Система розраховує профілі каркаса стелі та підлоговий і стельовий плінтуси & збіг з ручним розрахунком у метрах і штуках \\ \hline
    FR-10 & Система розраховує довжину кожного проводу за маршрутом і сумарну довжину із запасом & збіг з незалежним скриптом до 0,01~м \\ \hline
    FR-11 & Система будує розкладку листів ГКЛ у режимах «Авто», «Горизонтально», «Вертикально» й «Економія» з нумерацією плит, повторним використанням відрізків і списком залишків & нових листів не менше нижньої межі \eqref{eq:lowerbound}; сума статусів плит узгоджена \\ \hline
    FR-12 & Система зберігає проєкт у файл JSON та імпортує його з перевіркою й нормалізацією даних & завантажений проєкт відтворює сцену; хибні записи відкидаються \\ \hline
    FR-13 & Система дозволяє скасувати останні дії (до 60 кроків), підтримує гарячі клавіші збереження та скасування & після скасування стан сцени відновлюється \\ \hline
    FR-14 & Система дозволяє налаштувати відображення: сітку, підписи, розміри, лінійку дверей, кольори та зображення стін, дверей і картин & налаштування зберігаються в проєкті \\ \hline
\end{xltabular}
\end{spacing}

\customSubSubSection[sec:nfr]{Нефункціональні вимоги}{0.5em}

Нефункціональні вимоги структуровано за характеристиками моделі якості ISO/IEC~25010:2023~\citecustom{13} (\tabref{tab:nfr}).

\begin{spacing}{1.15}
\small\noindent
\begin{xltabular}{\textwidth}{| c | >{\raggedright\arraybackslash}p{3.0cm} | >{\raggedright\arraybackslash}X | >{\raggedright\arraybackslash}X |}
    \caption{Нефункціональні вимоги} \label{tab:nfr} \\
    \hline
    \textbf{ID} & \textbf{Характеристика} & \textbf{Вимога} & \textbf{Критерій} \\ \hline
    \endfirsthead

    \hline
    \textbf{ID} & \textbf{Характеристика} & \textbf{Вимога} & \textbf{Критерій} \\ \hline
    \endhead

    NFR-01 & Продуктивність & розкладка ГКЛ для сцени з двох кімнат виконується без відчутної затримки & час розрахунку $\le 0{,}5$~с у всіх режимах \\ \hline
    NFR-02 & Функціональна коректність & результати розрахунків збігаються з незалежним перерахунком & відхилення $\le 0{,}01$ (м$^2$, м) та збіг у штуках на контрольній сцені \\ \hline
    NFR-03 & Супроводжуваність та тестованість & алгоритм розкрою відокремлено від DOM і графічної бібліотеки & модуль запускається й перевіряється без застосунку \\ \hline
    NFR-04 & Автономність і розгортання & застосунок не потребує збірки, сервера застосунку та бази даних & запуск статичним HTTP-сервером \\ \hline
    NFR-05 & Приватність & дані проєкту не передаються на сервер & дані зберігаються лише у файлі користувача \\ \hline
    NFR-06 & Переносність & працює в сучасних браузерах із WebGL та ES-модулями & повна функціональність у Chromium; в інших~-- без перезапису файлу на місці \\ \hline
    NFR-07 & Надійність та цілісність даних & імпорт нормалізує дані, дії можна скасувати & відкидання недійсних записів; $\ge 60$ кроків скасування \\ \hline
    NFR-08 & Придатність до використання & основні сценарії виконуються без навчання & створення кімнати з дверима й розетками за $\le 10$~хв (планується оцінити юзабіліті-тестом) \\ \hline
\end{xltabular}
\end{spacing}

\customSubSubSection[sec:usecases]{Модель варіантів використання}{0.5em}

Основні варіанти використання та дійові особи показано на діаграмі UML~\citecustom{16} (рис.~\ref{fig:uc}).

\settikz[fig:uc]{0.82\textwidth}{assets/uc-diagram}{Діаграма варіантів використання вебзастосунку}

Ключовий сценарій «Розрахувати матеріали та побудувати розкрій» має такий опис.

\begin{itemize}[noitemsep, topsep=0.3em, leftmargin=*]
    \item \textbf{Передумова:} створено або імпортовано проєкт, у ньому є хоча б одна кімната.
    \item \textbf{Основний потік:} 1) користувач будує кімнати, двері, електрообладнання та проводи; 2) відкриває вкладку «Розрахунки»; 3) система обчислює площі, кількість листів, профілів, плінтусів і довжину проводів; 4) користувач обирає режим розкладки ГКЛ; 5) система будує розкладку, нумерує плити та формує список залишків; 6) користувач вмикає показ розкладки на 3D-моделі та зберігає проєкт.
    \item \textbf{Альтернативний потік A (хибні параметри):} якщо розмір листа неприпустимий, система виводить повідомлення й не будує розкладку.
    \item \textbf{Альтернативний потік B (немає поверхонь):} якщо не обрано ні стін, ні стелі, система повідомляє, що потрібно вибрати поверхні.
\end{itemize}

\customSubSection[sec:ch1-conclusions]{Висновки до розділу I}{1em}

\begin{enumerate}[noitemsep, topsep=0.3em, leftmargin=*, label=\arabic*.]
    \item Проблема полягає в ручному підрахунку довжин проводів і матеріалів за плоским планом, у відсутності маршруту проводу як просторового об'єкта в редакторах інтер'єру та в тому, що прості калькулятори не будують розкладку листів.
    \item Задачу розкрою листового матеріалу формалізовано як покриття чистої площі поверхонь гільйотинними розрізами з повторним використанням відрізків; виведено нижню межу кількості листів за площею.
    \item За результатами огляду трьох редакторів планування й двох онлайн-калькуляторів виявлено нішу: безкоштовний браузерний застосунок, що поєднує 3D-модель, трасування проводів і розкрій ГКЛ.
    \item Сформульовано 14 функціональних і 8 нефункціональних вимог із критеріями приймання, побудовано діаграму варіантів використання та описано ключовий сценарій.
\end{enumerate}
